Our preliminary findings reveal both the promise and current limitations of data-driven models in learning age-related motion signatures.
The ST-GCN age-signature detection network successfully segregates kinematic extremes --- Young versus Elderly --- confirming that motion clips contain discriminative biomechanical information related to age group.
However, the generated motions do not yet consistently reproduce statistically robust age-related differences across all selected clinical and kinematic metrics.
The LoRA-adapted Motion Diffusion Model successfully learns and imposes spatial kinematic constraints --- Elderly-conditioned clips produce measurably reduced hip Range-of-Motion relative to Young and Mid-age conditions --- yet spatiotemporal consistency remains unresolved, with generated stride variability massively overshooting clinical baselines.
We attribute this remaining failure to the capacity ceiling of low-rank, discrete-token adaptation: a rank-12 LoRA update can push a few high-leverage statistics such as hip RoM toward an age-specific mean, but cannot coherently re-shape stride timing, leaving the base MDM prior to dominate the temporal axis of the kinematic chain.
This gap highlights an important challenge for assistive robotics operating in eldercare and clinical environments: producing visually plausible age-conditioned motion is not sufficient for meaningful age-aware motion synthesis; motion prediction systems must also preserve biomechanical characteristics that affect anticipation, safety, and interaction.

The main bottleneck identified in this work is the limited size and quality of available age-diverse motion datasets.
Our experiments rely on clinical gait recordings collected from subjects across different age groups, but all available datasets are currently too small to support the training of deep-learning models, unavoidably resulting in model overfitting.

Another limitation is the current conditioning strategy.
While our proposed framework aims to condition generation on a continuous latent age representation extracted by the age-signature network, the implemented prototype uses discrete age-group conditions through LoRA adaptation.
This provides an efficient first step, but may not be expressive enough to capture gradual age-related changes or individual variability.

Future work will move beyond discrete age-group conditioning toward a continuous latent conditioning architecture (Fig.~\ref{fig:future_architecture}): a biomechanical style encoder is pretrained on AMASS via supervised contrastive learning to capture per-subject kinematic identity, an age head is fine-tuned on top of the frozen encoder using clinical age labels, and the MDM transformer is retrained end-to-end with conditioning formed by concatenating a CLIP text embedding $z_\text{text}$ and the biomechanical style embedding $z_\text{style}$.
At inference, a scalar age combined with a text prompt enables continuous interpolation across the age axis and personalization toward individual functional ability rather than coarse demographic bins.
